% Configurazione
\documentclass{article}

\usepackage{titling} % Required for inserting the subtitle
\usepackage{graphicx} % Required for inserting images
\usepackage{tabularx} % Per l'ambiente tabularx (tabelle)
\usepackage{calc} % Sempre per le tabelle
\usepackage[hidelinks]{hyperref} % Per i collegamenti ipertestuali, ad esempio sulla table of contents
\usepackage[italian]{babel} % Per la lingua italiana nelle scritte automatiche
\usepackage{xcolor} % Per colorare il testo
\usepackage{colortbl} % Per colorare le celle delle tabelle
\usepackage{lipsum} % Per generare lorem ipsum
\usepackage[normalem]{ulem} % Per sottolineare il testo
\usepackage{array} % Per la visualizzazione fluttuante di array di domande e risposte
\usepackage{ragged2e} % Pacchetto necessario per \justifying che giustifica il testo di tabelle

\newcommand{\ulhref}[2]{\href{#1}{\uline{#2}}} % Nuovo comando per sottolineare i link
\newcommand{\ulref}[1]{\uline{\ref{#1}}} % Nuovo comando per sottolineare i collegamenti a immagini
\setlength{\parindent}{0pt} % Rimuove il rientro automatico dei paragrafi

\graphicspath{ {immagini/} {../../../../shared/immagini/} }


% Variabili
\newcommand{\data}{30 Ottobre 2024}


% Struttura
\begin{document}

% Parte costante
\begin{minipage}{0.4\textwidth}
    \includegraphics[width=0.6\textwidth]{Logo Universita di Padova.png}
\end{minipage}
\begin{minipage}{0.55\textwidth}
    \textcolor{red}{\textbf{Università degli Studi di Padova}} \\
    \textcolor{red}{Laurea: Informatica} \\
    \textcolor{red}{Corso: Ingegneria del Software} \\
    \textcolor{red}{Anno Accademico: 2024/2025}
\end{minipage}

\begin{minipage}{0.4\textwidth}
    \includegraphics[width=0.6\textwidth]{Logo SWEg.png}
\end{minipage}
\begin{minipage}{0.55\textwidth}
    \textbf{Gruppo: SWEg Labs} \\
    Email: \textsf{gruppo.sweg@gmail.com}
\end{minipage}

\vspace{2cm}
{
\centering
\Huge\bfseries Verbale \par
\vspace{0.5cm}
\Large\bfseries \data\par
}
\newpage
\tableofcontents
\newpage

% Parte variabile
% Intestazione
\fancyhead[L]{1 \hspace{0.2cm} Informazioni generali} % Testo a sinistra

\pagenumbering{arabic} % Numerazione araba per il contenuto


\section{Informazioni generali}

\begin{itemize}
    \item \textbf{Tipo di riunione}: interna
    \item \textbf{Luogo}: meeting Discord
    \item \textbf{Data}: 06/02/2025
    \item \textbf{Ora inizio}: 17:30
    \item \textbf{Ora fine}: 19:00
    \item \textbf{Responsabile}: Giacomo Loat
    \item \textbf{Scriba}: Davide Verzotto
    \item \textbf{Partecipanti}:
    \begin{itemize}
        \item Federica Bolognini
        \item Michael Fantinato
        \item Giacomo Loat
        \item Filippo Righetto
        \item Riccardo Stefani
        \item Davide Verzotto
    \end{itemize}
\end{itemize}

\newpage
% Intestazione
\fancyhead[L]{2 \hspace{0.2cm} Ordine del giorno} % Testo a sinistra


\section{Ordine del giorno}

\begin{itemize}
    \item Analisi dei documenti richiesti per la \emph{RTB}\textsubscript{\textit{\textbf{G}}}, formulazione di strategie e suddivisione dei compiti per
    la loro stesura iniziale.
    \item Decidere in definitiva lo strumento per conservare e tracciare il \emph{backlog}\textsubscript{\textit{\textbf{G}}}.
    \item Decidere in definitiva lo strumenti per creare i \emph{diagrammi di Gantt}\textsubscript{\textit{\textbf{G}}} e i grafici sull'andalamento del 
    progetto.
    \item Discutere la \emph{best practice}\textsubscript{\textit{\textbf{G}}} di ingegneria del software che prevede di avere un 
    \emph{repository}\textsubscript{\textit{\textbf{G}}} unico per tutto il progetto, valutare il nostro caso specifico.
\end{itemize}

\newpage
% Intestazione
\fancyhead[L]{3 \hspace{0.2cm} Diario della riunione} % Testo a sinistra

\section{Diario della riunione}

\begin{itemize}
    \item Sono stati presentati al \emph{proponente} i nostri obiettivi di funzionalità per l'\emph{MVP}, ovvero: 
        \begin{itemize}
            \item Salvataggio e recupero dello storico di sessione;
            \item Proporre domande per iniziare la conversazione;
            \item Proporre domande per proseguire la conversazione;
            \item Visualizzazione del badge dell’esito dell’aggiornamento automatico;
            \item Migliorare l’aggiornamento automatico per gestire l’eliminazione di file;
            \item Migliorare l’aggiornamento automatico dei file che subiscono modifiche.
        \end{itemize}
    \item Sono state presentate al \emph{proponente} le funzionalità che verranno trattate per ultime, nel caso in cui avessimo
    ancora del tempo a disposizione, ovvero:
        \begin{itemize}
            \item Visualizzazione dei file da cui il bot ha preso la risposta. Inoltre, per questa funzionalità è stato scelto di
            presentare il link al file sulla piattaforma;
            \item Integrazione delle \emph{API}\textsubscript{\textbf{\textit{G}}} verso
            \emph{Telegram}\textsubscript{\textbf{\textit{G}}} e \emph{Slack}\textsubscript{\textbf{\textit{G}}}.
        \end{itemize}

    \item Abbiamo fornito al \emph{proponente} un aggiornamento sullo stato attuale della progettazione, ovvero: 
        \begin{itemize}
            \item Abbiamo riferito che stiamo applicando l'\emph{architettura esagonale}\textsubscript{\textbf{\textit{G}}};
            \item Abbiamo riferito che abbiamo scelto un sistema in cui una sessione è il periodo di tempo in cui l'app è accesa.
            Inoltre abbiamo chiesto approvazione per l'utilizzo del \emph{Singleton}\textsubscript{\textbf{\textit{G}}} per modellare
            la sessione corrente.
        \end{itemize}
    \item Abbiamo chiesto consigli al \emph{proponente} riguardanti la progettazione;
    \item Abbiamo discusso a riguardo del concetto di sessione e abbiamo accettato la proposta di rimuovere tale suddivisione, e piuttosto
    recuperare lo storico dei messaggi in modo amalgamato;
    \item Il \emph{proponente} ci ha chiarito che il software che ci è richiesto verrà poi ulteriormente sviluppato internamente
    all'azienda, e quindi ci ha suggerito di svolgere una progettazione dell'\emph{MVP} in funzione delle necessità future di sviluppo, 
    in modo da facilitare il futuro lavoro dei programmatori di \emph{AzzurroDigitale}, piuttosto che ragionare come se il prodotto
    fosse destinato a non essere più sviluppato dopo la nostra consegna;
    \item Abbiamo esposto al \emph{proponente} i vari test che dovremmo implementare, ovvero:
        \begin{itemize}
            \item Abbiamo riferito che abbiamo scelto i \emph{framework}\textsubscript{\textbf{\textit{G}}}
            \emph{PyTest}\textsubscript{\textbf{\textit{G}}} e \emph{Jasmine}\textsubscript{\textbf{\textit{G}}}, e abbiamo riferito
            che prepareremo l'ambiente di test a livello di struttura di cartelle;
            \item Abbiamo riferito che per i test di unità a breve sceglieremo i metodi più importanti da testare.
        \end{itemize}
    \item Abbiamo effettuato un resoconto al \emph{proponente} sullo stato di avanzamento relativo alla documentazione. In particolare,
    abbiamo discusso a riguardo del \emph{Manuale Utente}\textsubscript{\textbf{\textit{G}}}, sottolineando l'importanza di fornire
    una guida su "Cosa chiedere e come chiederlo", e dunque di scrivere con particolare cura tale sezione;
    \item Abbiamo confermato lo scheduling dei prossimi incontri e la durata delle sprint pari a 2 settimane.


\end{itemize}
\newpage
% Intestazione
\fancyhead[L]{5 \hspace{0.2cm} Todo} % Testo a sinistra

\section{Todo}

Durante la riunione sono emersi i seguenti task da svolgere:

\vspace{0.5cm}

\begin{table}[htbp]
\centering
\rowcolors{2}{lightgray}{white}
\begin{tabular}{|c|c|p{0.25\textwidth}|p{0.4\textwidth}|}
    \hline
    \rowcolor[gray]{0.75}
    \textbf{Codice} & \textbf{Dalla decisione} & \textbf{Assegnatario} & \textbf{Task Todo} \\
    \hline
    BUD-320 & VI 30.4 & Riccardo Stefani & Architettura della generazione di una risposta.\\
    \hline
    BUD-321 & VI 30.4 & Riccardo Stefani & Architettura dell’aggiornamento automatico del \emph{database vettoriale}\textsubscript{\textit{\textbf{G}}}.\\
    \hline
    BUD-322 & VI 30.4 & Michael Fantinato & Architettura dell’inizializzazione, refresh e scroll di \emph{Angular}.\\
    \hline 
    BUD-323 & VI 30.4 & Michael Fantinato & Architettura dell’aggiornamento del rendering grafico durante la chat.\\
    \hline 
    BUD-324 & VI 30.4 & Michael Fantinato & Architettura frontend dell’aggiornamento del badge di segnalazione dell’esito dell’ultimo aggiornamento automatico.\\
    \hline 
    BUD-325 & VI 30.4 & Michael Fantinato & Architettura backend dell’aggiornamento del badge di segnalazione dell’esito dell’ultimo aggiornamento automatico.\\
    \hline 
    BUD-326 & VI 30.4 & Giacomo Loat & Architettura del salvataggio dei messaggi nello storico.\\
    \hline 
    BUD-327 & VI 30.4 & Giacomo Loat & Architettura del recupero dei messaggi dallo storico.\\
    \hline 
    BUD-328 & VI 30.4 & Riccardo Stefani & Architettura della generazione di domande per proseguire la conversazione.\\
    \hline 
    BUD-336 & VI 30.7 & Michael Fantinato & Modellare in \emph{Angular} il pulsante di copia messaggio e il pulsante di copia \emph{snippet} di codice. \\
    \hline
    BUD-319 & VI 30.8 & Federica Bolognini & Redigere il verbale interno del 24/02/2025. \\
    \hline

\end{tabular}
\end{table}

\end{document}
